\subsection{An arithmetic cross-cup witness with trivial word and CM actions}
\label{sec:cross-cup-residual-inertia}

\begin{proposition}\label{prop:cross-cup-residual-inertia}
Choose pro-two elements \(\tau_5\in I_5\) and
\(\tau_{29}\in I_{29}\) of arithmetic inertia with
\(c_5(\tau_5)=c_{29}(\tau_{29})=1\).
Then
\[
                         w=[\tau_5,\tau_{29}]
\]
belongs to \(V_D\cap H\). In fact every rational Kummer
coordinate, every free pro-\(p\) action, and both
complete corrected CM frames vanish on \(w\).
The character \(\kappa_{5,29}\) of
\[
 \mathbf Q(\sqrt5,\sqrt{29},\sqrt{7+2\sqrt5})
\]
has value one on \(w\).
\end{proposition}

\begin{proof}
Cyclotomic inertia at an odd prime \(\ell\) is supported
only at the coefficient prime \(\ell\), and its projection
onto the quadratic character \(c_\ell\) is surjective.
Taking the pro-two power of a suitable inertia element
therefore supplies the stated \(\tau_\ell\).

Put \(u=\lambda(\tau_5)\) and \(v=\lambda(\tau_{29})\).
Both are one at the coefficient prime two, while \(u\)
is one away from five and \(v\) is one away from twenty-nine.
For any rational Kummer cocycle \(\rho_q\), the crossed
homomorphism identity gives, componentwise,
\begin{equation}
 \rho_q([\tau_5,\tau_{29}])
       =(1-v)\rho_q(\tau_5)+(u-1)\rho_q(\tau_{29}).        \label{eq:cross-cup-inertia-cocycle}
\end{equation}
At two both coefficients vanish.
At five the only possible term is
\((u_5-1)\rho_{q,5}(\tau_{29})\), and this is zero:
on the pro-two cyclic group generated by \(\tau_{29}\),
the fifth multiplier is one, so \(\rho_{q,5}\) is a
continuous homomorphism to the pro-five additive group.
At twenty-nine the symmetric argument applies.
At every other prime both coefficients vanish.
Thus \eqref{eq:cross-cup-inertia-cocycle} is zero in the full adelic target, for every
rational \(q\), including \(q=2\). The commutator also
has full multiplier one. Therefore \(w\in D\cap H\).

Arithmetic inertia at \(\ell\) acts trivially in the
free pro-\(p\) action for \(p\ne\ell\), by
tangential unramifiedness
\cite[\S8.5, p.~151]{HainMatsumoto2003}.
For \(p=5\), the second commutator entry has trivial
action; for \(p=29\), the first does; at every other
prime both do. Hence \(w\in P_D\cap H\).
The cubic defect vanishes on \(D\cap H\), so
\(w\in T_D\cap H\).

Both inertia elements fix \(\mathbf Q(i)\) and
\(\mathbf Q(\sqrt{-3})\), whose finite ramification is
confined to two and three.
Their CM Tate actions therefore belong to the
respective commutative CM unit groups.
Both Tate commutators are the identity.
On \(D\cap H\), the established signed comparisons
identify the corrected frames with these Tate actions;
the central cubic correction is trivial on \(D\cap H\).
Thus \(T(w)=W(w)=1\), in every coefficient component,
and in particular \(w\in V_D\cap H\).

In the dihedral field defining \(\kappa_{5,29}\),
the images of \(\tau_5\) and \(\tau_{29}\) have respective
biquadratic coordinates \((1,0)\) and \((0,1)\):
the other quadratic character is unramified at that
inertia prime.
For its central extension class \(c_5c_{29}\), the
commutator of such two lifts is the central generator,
independently of the choices of lifts.
Thus \(\kappa_{5,29}(w)=1\).
\end{proof}

The proposition proves arithmetic nonvanishing on the
residual subgroup $V_D$ itself. It does not assert
that \(\kappa_{5,29}\) extends to a full \(G\)-invariant
character. That is a separate cohomological lifting
question.
